\subsection{Повороты и балансировка}

\Def Балансировкой AVL-дерева будем называть процесс, который из
дерева поиска, в узле которого разность высот поддеревьев равна
двум, переподвешивает детей и внуков так, чтобы выполнялось
свойство AVL-дерева.

Пусть $h(v)$ - высота поддерева с корнем в $v$, $\delta(v) = h(L) - h(R)$, где $h(L)$ и $h(R)$ - высоты поддеревьев вершины $v$. 

\subsubsection*{Малый левый поворот}

\begin{itemize}
    \item Используется, когда $\delta(a) = -2$ и $\delta(b) \leq 0$
    \drawsome{pix/avl_left_rotate.png}
\end{itemize}

\subsubsection*{Большой левый поворот}

\begin{itemize}
    \item Используется, когда $\delta(a) = -2$ и $\delta(b) > 0$
    \drawsome{pix/avl_big_left_rotate.png}
\end{itemize}

\subsubsection*{Малый правый поворот}

\begin{itemize}
    \item Используется, когда $\delta(a) = 2$ и $\delta(b) \geq 0$
    \drawsome{pix/avl_right_rotate.png}
\end{itemize}

\subsubsection*{Большой правый поворот}

\begin{itemize}
    \item Используется, когда $\delta(a) = 2$ и $\delta(b) < 0$
    \drawsome{pix/avl_big_right_rotate.png}
\end{itemize}

\Statement Все вращения требуют $O(1)$ операций. 

\subsubsection{Корректность}

\drawsomemedium{pix/avl_left_rotate.jpg}

\Proof 
Рассмотрим малый левый поворот. Так как все поддерево - это чей-то правый или левый ребенок, то не имеет значения, за какую вершину этого поддерева мы подвесим поддерево к родителю.
\begin{itemize}
    \item R - правое поддерево b. До поворота оно тоже было правым поддеревом b, так что эта связь корректна.
    \item a - левый ребенок b, поэтому b > a. Таким образом, связь между a и b в дереве после поворота корректна.
    \item Q лежало в левом поддереве b, т.е. все значения Q меньше или равны b. Значит, связь Q-b корректна. Q - правый ребенок a, поэтому то, что Q лежит в правом поддереве a, тоже корректно.
    \item P - левое поддерево a. До поворота оно тоже было левым поддеревом a, так что эта связь корректна. R лежало в правом поддереве b, а P лежало в левом. Поэтому R > P, так что то, что P лежит в левом поддереве b, - корректно.
\end{itemize}

Для правого поворота аналогично.

Большой поворот - это комбинация 2х малых поворотов, поэтому большие повороты
также корректны.
\Endproof

\subsection{Основные операции}

\subsubsection*{Поиск элемента.}

Поиск в дереве не меняется.

\subsubsection*{Вставка элемента.}

Необходимо добавить вершину со значением $x$. 

\Alg
\begin{enumerate}
    \item Вставляем элемент как в обычном двоичном дереве.
    Вершина со значением $x$ добавлена, как лист. Баланс в дереве мог нарушиться только в вершинах, которые находятся по пути от корня до добавленной вершины.
    \item Поднимаемся от вершины к корню. Проверяем сбалансированность. Если разность высот поддеревьев равна 2 – выполняем нужное вращение. (Считаем, что разность высот поддеревьев не может стать больше 2. Это не просили доказывать, но проверит факт просто. Нужно посмотреть, как меняется баланс в дереве при вращении, добавлении, удалении)
\end{enumerate}

\textit{Асимптотика.} Добавление за вершины как лист: $O(log(n))$. Так как высота AVL-дерева $O(log(n))$, а вращение происходит за $O(1)$, то суммарное время добавление элемента в дерево $O(log(n))$. 

\subsubsection*{Удаление элемента.}

Необходимо удалить вершину со значением $x$. 

\Alg
\begin{enumerate}
    \item Удаляем элемент как в двоичном дереве поиска
    Баланс в дереве мог нарушиться только в вершинах, которые находятся по пути от корня до удалённой вершины.
    \item Поднимаемся от места удалённой вершины к корню. (в и по вершине, которая встало на место удалённой. (если такая есть)) Проверяем сбалансированность. Если разность высот поддеревьев равна 2 – выполняем нужное вращение.
\end{enumerate}

\textit{Асимптотика.} Нахождение и удаление вершины из дерева: $O(log(n))$. Так как высота AVL-дерева $O(log(n))$, а вращение происходит за $O(1)$, то суммарное время удаления элемента в дереве $O(log(n))$. 

\pagebreak